\section{One-sided bounds on windows}\label{sec:criterion}

This section proves the analytic criterion of the paper. An upper bound $G_+(x)\le Cx^\theta$ for a positive combination of smoothed prime sums, known for $x$ in a sparse sequence of windows $Y\le x\le Y^\lambda$, excludes zeros of $\zeta$ and of the corresponding Dirichlet $L$-functions from the half-plane $\operatorname{Re}s>\theta$. The proof is a quantitative form of Landau's theorem on Dirichlet integrals of functions of one sign~\cite{landau}.

\subsection{The smoothed prime sums}\label{sec:functions}

Let $\Lambda$ denote the von Mangoldt function and $\gamma$ Euler's constant. For $x\ge1$ put
\begin{equation}\label{G}
G(x)=\sum_{n\le x}\Lambda(n)\Big(\frac xn-1\Big)-x\log x+(1+\gamma)x ,
\end{equation}
and for a nontrivial Dirichlet character $\chi$ put
\begin{equation}\label{Gchi}
G_\chi(x)=\sum_{n\le x}\Lambda(n)\operatorname{Re}\chi(n)\Big(\frac xn-1\Big)+\operatorname{Re}\frac{L'(1,\chi)}{L(1,\chi)}\,x .
\end{equation}
For $\operatorname{Re}s>1$ and every integer $n\ge1$,
\begin{gather*}
\int_n^\infty\Big(\frac xn-1\Big)x^{-s-1}\,dx=\frac{n^{-s}}{s(s-1)},\\
\int_1^\infty x\log x\;x^{-s-1}\,dx=\frac1{(s-1)^2},\qquad
\int_1^\infty x\;x^{-s-1}\,dx=\frac1{s-1}.
\end{gather*}
Summation over $n$ gives, for $\operatorname{Re}s>1$,
\begin{equation}\label{F}
\int_1^\infty G(x)x^{-s-1}\,dx=F(s),\qquad \int_1^\infty G_\chi(x)x^{-s-1}\,dx=\tfrac12\big(F_\chi(s)+F_{\bar\chi}(s)\big),
\end{equation}
where
\begin{equation}\label{Fchi}
F(s)=-\frac{\zeta'(s)}{\zeta(s)}\,\frac1{s(s-1)}-\frac1{(s-1)^2}+\frac{1+\gamma}{s-1},\qquad
F_\chi(s)=-\frac{L'(s,\chi)}{L(s,\chi)}\,\frac1{s(s-1)}+\frac{L'(1,\chi)}{L(1,\chi)}\,\frac1{s-1}.
\end{equation}
The expansions $-\zeta'(s)/\zeta(s)=(s-1)^{-1}-\gamma+O(s-1)$ and $1/(s(s-1))=(s-1)^{-1}-1+O(s-1)$ show that $F$ is analytic at $s=1$. The same computation, together with $L(1,\chi)\ne0$, shows that $F_\chi$ is analytic at $s=1$. Both functions are meromorphic on $\mathbb C$. In the half-plane $\operatorname{Re}s>0$ their poles are simple and lie at the zeros $\rho$ of $\zeta$ and of $L(\cdot,\chi)$, with residues
\begin{equation}\label{res}
\Res_{s=\rho}F=-\frac{m_\zeta(\rho)}{\rho(\rho-1)},\qquad \Res_{s=\rho}F_\chi=-\frac{m_\chi(\rho)}{\rho(\rho-1)},
\end{equation}
where $m_\zeta(\rho)$ and $m_\chi(\rho)$ denote the orders of vanishing of $\zeta$ and of $L(\cdot,\chi)$ at $\rho$.

We use four classical facts~\cite{davenport}. The functions $\zeta$ and $L(\cdot,\chi)$ do not vanish on the half-plane $\operatorname{Re}s\ge1$, and $\zeta$ does not vanish on the interval $0<s<1$. Mertens' estimate $\sum_{n\le x}\Lambda(n)/n=\log x+O(1)$ and Chebyshev's estimate $\sum_{n\le x}\Lambda(n)=O(x)$ hold, and for nontrivial $\chi$ the partial sums $\sum_{n\le x}\Lambda(n)\chi(n)/n$ are bounded. These estimates give a constant $B$ with
\begin{equation}\label{triv}
\lvert G(x)\rvert\le Bx,\qquad \lvert G_\chi(x)\rvert\le Bx\qquad(x\ge1).
\end{equation}
The function $G$ is continuous, because the term of index $n$ in \eqref{G} vanishes at $x=n$, and away from the integers $G'(x)=\sum_{n\le x}\Lambda(n)/n-\log x+\gamma$ is bounded. Hence there is a constant $K$ with
\begin{equation}\label{lip}
\lvert G(x)-G(y)\rvert\le K\lvert x-y\rvert\qquad(x,y\ge1),
\end{equation}
and the same argument gives \eqref{lip} for $G_\chi$.

\subsection{Three lemmas}

\begin{lemma}\label{tails}
Let $k\ge0$ be an integer, let $U\ge0$ and $a>0$. If $b>a$, then $\int_0^Ue^{-au}\frac{u^k}{k!}\,du\le e^{(b-a)U}b^{-k-1}$. If $0<b<a$, then $\int_U^\infty e^{-au}\frac{u^k}{k!}\,du\le e^{-(a-b)U}b^{-k-1}$.
\end{lemma}

\begin{proof}
For $b>a$ and $0\le u\le U$ we have $e^{-au}=e^{-bu}e^{(b-a)u}\le e^{(b-a)U}e^{-bu}$. For $0<b<a$ and $u\ge U$ we have $e^{-au}=e^{-bu}e^{-(a-b)u}\le e^{-(a-b)U}e^{-bu}$. In both cases the identity $\int_0^\infty e^{-bu}\frac{u^k}{k!}\,du=b^{-k-1}$ completes the proof.
\end{proof}

\begin{lemma}\label{rot}
Let $0<\varphi<\pi/2$ and let $\alpha$ be real. Every set of $\lceil2\pi/\varphi\rceil$ consecutive integers contains an integer $j$ with $\cos(\alpha+j\varphi)\ge\frac12$.
\end{lemma}

\begin{proof}
The set of angles with cosine at least $\frac12$ is the union of the arcs $[-\pi/3,\pi/3]+2\pi\mathbb Z$, of length $2\pi/3$, separated by gaps of length $4\pi/3$. The angles $\alpha+j\varphi$ advance in steps of length $\varphi<2\pi/3$, so they cannot jump across an arc. A sequence of angles that starts in a gap leaves it after at most $\lceil4\pi/(3\varphi)\rceil$ steps, and its first angle outside the gap lies in an arc. The count $\lceil4\pi/(3\varphi)\rceil+1$ is at most $\lceil2\pi/\varphi\rceil$ because $2\pi/\varphi-4\pi/(3\varphi)=2\pi/(3\varphi)>1$.
\end{proof}

\begin{lemma}\label{near}
Let $P\subset\mathbb C$ be a closed discrete set that is stable under complex conjugation. Let $\theta<1$, let $\rho_0=\beta_0+i\gamma_0\in P$ with $\theta<\beta_0<1$, and put $\theta'=(\theta+\beta_0)/2$. Let $\sigma_1$ be a real number with $\sigma_1>1+(\lvert\rho_0\rvert^2+\gamma_0^2)/(\beta_0-\theta)$, and let $r$ be the distance from $\sigma_1$ to $P$. Then $r<\sigma_1-\theta'$, and every $p\in P$ with $\lvert\sigma_1-p\rvert=r$ satisfies $\operatorname{Re}p>\theta'$. For all such $\sigma_1$ outside a countable set, the points of $P$ at distance $r$ from $\sigma_1$ form a single real point or a single pair of complex conjugate points.
\end{lemma}

\begin{proof}
For $p\in P$ with $\operatorname{Re}p\le\theta'$,
\[
\lvert\sigma_1-p\rvert^2-\lvert\sigma_1-\rho_0\rvert^2=2\sigma_1(\beta_0-\operatorname{Re}p)+\lvert p\rvert^2-\lvert\rho_0\rvert^2\ge\sigma_1(\beta_0-\theta)-\lvert\rho_0\rvert^2>0,
\]
so $p$ lies farther from $\sigma_1$ than $\rho_0$ does. Next,
\[
(\sigma_1-\theta')^2-\lvert\sigma_1-\rho_0\rvert^2=(\beta_0-\theta')(2\sigma_1-\beta_0-\theta')-\gamma_0^2\ge(\beta_0-\theta)(\sigma_1-1)-\gamma_0^2>0,
\]
so $r\le\lvert\sigma_1-\rho_0\rvert<\sigma_1-\theta'$. Finally, two points $p,p'\in P$ are equidistant from $\sigma_1$ exactly when $2\sigma_1(\operatorname{Re}p'-\operatorname{Re}p)=\lvert p'\rvert^2-\lvert p\rvert^2$. If $\operatorname{Re}p\ne\operatorname{Re}p'$, this equation holds for one value of $\sigma_1$, and if $\operatorname{Re}p=\operatorname{Re}p'$, it forces $p'=p$ or $p'=\bar p$. The set $P$ is countable, so the exceptional values of $\sigma_1$ form a countable set.
\end{proof}

\subsection{The window theorem}\label{sec:win}

\begin{theorem}\label{engine}
Let $H\colon[1,\infty)\to\mathbb R$ be measurable with $\lvert H(x)\rvert\le Bx$ for a constant $B$, and suppose that the function $A(s)=\int_1^\infty H(x)x^{-s-1}\,dx$, defined for $\operatorname{Re}s>1$, extends to a meromorphic function on $\mathbb C$. Let $0\le\theta<1$. Suppose that every pole $p$ of $A$ with $\operatorname{Re}p>\theta$ is simple and satisfies $\operatorname{Re}p<1$, and that the residue of $A$ at every real pole $p>\theta$ is positive. Suppose that there are constants $\lambda>1$ and $C$ such that, for arbitrarily large $Y$, the inequality $H(x)\le Cx^\theta$ holds for every $x$ with $Y\le x\le Y^\lambda$. Then $A$ does not have a pole in the half-plane $\operatorname{Re}s>\theta$.
\end{theorem}

The proof compares two evaluations of the Taylor coefficients $(-1)^kA^{(k)}(\sigma_1)/k!$ at a real point $\sigma_1>1$. The nearest pole of $A$ forces these coefficients to be large for a positive proportion of $k$. The window bound forces them to be small for every $k$ in a range proportional to $\log Y$, because the kernel $e^{-\sigma_1u}u^k/k!$ concentrates near $u=k/(\sigma_1-1)$ with relative width of order $k^{-1/2}$.

\begin{proof}
Suppose that $A$ has a pole $\rho_0=\beta_0+i\gamma_0$ with $\beta_0>\theta$, and let $P$ be the set of poles of $A$. Since $H$ is real, $A(\bar s)=\overline{A(s)}$, so $P$ is stable under conjugation and the residues at conjugate poles are conjugate.

\emph{Step 1. The constants.} Fix $\mu,\nu\in(0,1)$ with $\mu/\log(1+\mu)<\lambda^{1/3}$ and $-\log(1-\nu)/\nu<\lambda^{1/3}$. For real numbers $a>0$ and $\eta\ge0$ define
\[
\xi_-=a\,\frac{(1+\mu)e^{\eta}-1}{\log(1+\mu)},\qquad \xi_+=\frac{a\,\nu\lambda}{\eta-\log(1-\nu)} .
\]
At $\eta=0$ we have $\xi_-/a<\lambda^{1/3}<\lambda^{2/3}<\xi_+/a$. Both ratios are continuous in $\eta$ and independent of $a$, so there is $\eta_0>0$ with $\xi_-<\xi_+$ whenever $0\le\eta\le\eta_0$.

\emph{Step 2. The base point.} Put $\theta'=(\theta+\beta_0)/2$. Lemma~\ref{near} gives a real number $\sigma_1$, as large as we please, such that $r=\operatorname{dist}(\sigma_1,P)$ satisfies $r<\sigma_1-\theta'$, and such that the poles nearest to $\sigma_1$ form either one real pole $p_1$ or one conjugate pair $p_1,\bar p_1$ with $\operatorname{Im}p_1>0$. In both cases $\theta'<\operatorname{Re}p_1<1$. Put $a=\sigma_1-1$ and $\eta=\log(r/a)$. Then $r>a$, and $\eta\le\log\big((\sigma_1-\theta')/(\sigma_1-1)\big)$, which is at most $\eta_0$ once $\sigma_1$ is large. We fix such a $\sigma_1$, we fix $\xi_1<\xi_2$ in the interval $(\xi_-,\xi_+)$, and we put $\kappa=(\sigma_1-\theta)/(\sigma_1-\theta')>1$.

\emph{Step 3. The Taylor coefficients from the poles.} For an integer $k\ge0$ put
\[
T_k=\int_0^\infty H(e^u)\,e^{-\sigma_1u}\,\frac{u^k}{k!}\,du=\frac{(-1)^k}{k!}A^{(k)}(\sigma_1).
\]
Let $R_1$ be the residue of $A$ at $p_1$. Subtracting the principal parts of $A$ at the poles nearest to $\sigma_1$ leaves a function $Q$ that is analytic on a disc $\lvert s-\sigma_1\rvert\le r''$ with $r''>r$, and Cauchy's estimate gives $\lvert Q^{(k)}(\sigma_1)\rvert/k!\le M(r'')^{-k}$ with $M=\max_{\lvert s-\sigma_1\rvert=r''}\lvert Q(s)\rvert$. Since $\frac{(-1)^k}{k!}\frac{d^k}{ds^k}(s-p)^{-1}=(s-p)^{-k-1}$, a real nearest pole gives
\[
T_k=R_1\,r^{-k-1}+\frac{(-1)^k}{k!}Q^{(k)}(\sigma_1),\qquad R_1>0 .
\]
For a conjugate pair write $\sigma_1-p_1=re^{-i\varphi}$ with $0<\varphi<\pi/2$ and $R_1=\lvert R_1\rvert e^{i\alpha}$. Then
\[
T_k=2\lvert R_1\rvert\,r^{-k-1}\cos\big(\alpha+(k+1)\varphi\big)+\frac{(-1)^k}{k!}Q^{(k)}(\sigma_1).
\]
By Lemma~\ref{rot} there is an integer $K_0$ such that every set of $K_0$ consecutive integers contains a $k$ with $\cos(\alpha+(k+1)\varphi)\ge\frac12$. Call such $k$ good, and in the real case call every $k$ good. Since $r''>r$, there is $k_0$ such that every good $k\ge k_0$ satisfies
\begin{equation}\label{lower}
T_k\ge\tfrac12\lvert R_1\rvert\,r^{-k-1}.
\end{equation}

\emph{Step 4. The Taylor coefficients from the window.} Let $Y$ be one of the values in the hypothesis and put $L=\log Y$. We split the integral defining $T_k$ at $u=L$ and at $u=\lambda L$. On the middle range $H(e^u)\le Ce^{\theta u}$, hence
\[
\int_L^{\lambda L}H(e^u)e^{-\sigma_1u}\frac{u^k}{k!}\,du\le C(\sigma_1-\theta)^{-k-1}\le C\kappa^{-k-1}r^{-k-1},
\]
where the last inequality uses $r<\sigma_1-\theta'$. On the outer ranges $\lvert H(e^u)\rvert e^{-\sigma_1u}\le Be^{-au}$. Lemma~\ref{tails} with $b=(1+\mu)r$ on $[0,L]$ and with $b=(1-\nu)a$ on $[\lambda L,\infty)$ gives
\begin{align*}
\int_0^L Be^{-au}\frac{u^k}{k!}\,du&\le B\,r^{-k-1}\exp\Big(\log(1+\mu)\big(\xi_-L-(k+1)\big)\Big),\\
\int_{\lambda L}^\infty Be^{-au}\frac{u^k}{k!}\,du&\le B\,r^{-k-1}\exp\Big(\big(\eta-\log(1-\nu)\big)\big((k+1)-\xi_+L\big)\Big),
\end{align*}
since $(1+\mu)r-a=\xi_-\log(1+\mu)$ and $r/((1-\nu)a)=e^{\eta}/(1-\nu)$. For every integer $k$ with $\xi_1L\le k+1\le\xi_2L$ both exponents are at most $-\delta L$, where
\[
\delta=\min\big\{(\xi_1-\xi_-)\log(1+\mu),\,(\xi_+-\xi_2)(\eta-\log(1-\nu))\big\}>0 .
\]
For every such $k$ we therefore have
\begin{equation}\label{upper}
T_k\le r^{-k-1}\big(2Be^{-\delta L}+C\kappa^{-k-1}\big).
\end{equation}

\emph{Step 5. The contradiction.} The interval $[\xi_1L,\xi_2L]$ has length $(\xi_2-\xi_1)L$, so for large $L$ it contains $K_0$ consecutive values of $k+1$ with $k\ge k_0$, and one of these $k$ is good. For that $k$, the inequalities \eqref{lower} and \eqref{upper} give $\frac12\lvert R_1\rvert\le2Be^{-\delta L}+C\kappa^{-\xi_1L}$. The right side tends to zero as $L\to\infty$ along the admissible values of $Y$, while $\lvert R_1\rvert>0$. Hence $A$ does not have a pole with real part greater than $\theta$.
\end{proof}

\begin{corollary}\label{corA}
Let $\chi_1,\dots,\chi_N$ be nontrivial Dirichlet characters, let $w_1,\dots,w_N$ be positive real numbers, and put $G_+=G+\sum_{i=1}^Nw_iG_{\chi_i}$ with $G$ and $G_\chi$ as in \eqref{G} and \eqref{Gchi}. Let $0\le\theta<1$, and suppose that there are constants $\lambda>1$ and $C$ such that, for arbitrarily large $Y$, the inequality $G_+(x)\le Cx^\theta$ holds for every $x$ with $Y\le x\le Y^\lambda$. Then $\zeta(s)\ne0$ and $L(s,\chi_i)\ne0$ for every $i$ and every $s$ with $\operatorname{Re}s>\theta$.
\end{corollary}

\begin{proof}
By \eqref{F}, \eqref{Fchi} and \eqref{triv}, the function $H=G_+$ satisfies $\lvert H(x)\rvert\le(1+\sum_iw_i)Bx$, and its transform $A=F+\frac12\sum_iw_i(F_{\chi_i}+F_{\bar\chi_i})$ is meromorphic on $\mathbb C$. By \eqref{res}, the poles of $A$ in $\operatorname{Re}s>\theta$ are simple, they lie at the zeros $\rho$ of $\zeta$, of the $L(\cdot,\chi_i)$ and of the $L(\cdot,\bar\chi_i)$, and
\[
\Res_{s=\rho}A=-\frac{c_\rho}{\rho(\rho-1)},\qquad c_\rho=m_\zeta(\rho)+\tfrac12\sum_{i=1}^Nw_i\big(m_{\chi_i}(\rho)+m_{\bar\chi_i}(\rho)\big)\ge0 .
\]
These zeros satisfy $\operatorname{Re}\rho<1$. At a real zero $\rho\in(\theta,1)$ the residue is positive because $\rho(\rho-1)<0$. Theorem~\ref{engine} shows that $A$ does not have a pole with $\operatorname{Re}s>\theta$, so $c_\rho=0$ for every $\rho$ in that half-plane. Each term of $c_\rho$ is nonnegative and each $w_i$ is positive, so each multiplicity vanishes.
\end{proof}

If $E$ is an abelian number field with group $X$ of primitive Dirichlet characters, then $\zeta_E(s)=\prod_{\chi\in X}L(s,\chi)$, and the weights $w_\chi=1$ for the nontrivial $\chi\in X$ make $G_+$ the smoothed sum over the prime ideals of $E$. Part~II proves a bound of this kind for a positive combination of characters attached to one CM field, and Corollary~\ref{corA} then treats all of these characters at once.
